\newif\ifglobal

%%%%%%%%%%%%%%%%%%%%%%%
%%% Compile to view %%%
%%%%%%%%%%%%%%%%%%%%%%%

%%%%%%%%%%%%%%%%%%%%%%%%%%%%%%%%%%%%%%%%%%%%%%%%%%%
%%% Readme:                                     %%%
%%% https://www.overleaf.com/read/bwdjvfjksvjy/ %%%
%%%%%%%%%%%%%%%%%%%%%%%%%%%%%%%%%%%%%%%%%%%%%%%%%%%

% >>> M?: marks a module setting number or important notice

% >>> M4: Set {./relative-path} to external files for this module <<<
\def \modulefiles {./34-RNG}
% To add an external file, use:
% (Replace '---' with actual filename)
% '\input{\modulefiles/tables/---.tex}' for tables
% '\includegraphics{\modulefiles/figures/---}' for figures
% '\subfile{\modulefiles/---}' for register files

% >>> M1: This module is in global project? <<<
% 'YES' - uncomment; 'NO' - comment out
\globaltrue

\ifglobal
    \documentclass[main\_\_EN.tex]{subfiles}

\else
    \documentclass[openany, 10pt]{book}

    % >>> M2: In ./chip-spec-settings/preamble-trm-module, also find
    %         settings for visibility of todo notes, watermark <<<
    \usepackage{preamble-trm-module}

    % >>> M3: Temporary labels in English,
    %         you can add your own labels there <<<
    \myexternaldocument{temp-labels-en}

\fi %\ifglobal


\setcounter{secnumdepth}{4}




% >>> M5: If this module needs any updates in
% './chip-spec-settings/preamble-trm-module.sty'
% (include additional package, etc.), contact your module PM
% or create the file './chip-spec-settings/changelog.tex' make the
% necessary updates yourself and document them in this file <<<

\raggedright


\begin{document}


% Sets language into which pre-defined bits of text are translated
\selectlanguage{english}

% Do not delete this block
\ifglobal
% Add the following front matter if
% this module is outside of global project
\else
    \InputIfFileExists{readme}{}{}
    {\let\clearpage\relax \listoftodos}
    {\let\clearpage\relax \listofdonetodos}
    \tableofcontents
    %\listoftables
    %\listoffigures
\fi


%%%%%%%%%%%%%%%%%%%%%%%%%
%%% TRM Module Begins %%%
%%%%%%%%%%%%%%%%%%%%%%%%%

\hypertarget{rng}{}
\chapter{Random Number Generator (RNG)}
\label{mod:rng}


\section{Introduction}

The \chipname{} contains a true random number generator, which generates 32-bit random numbers that can be used for cryptographical operations, among other things.

%%%%%%%%%%%%%%%%%%%%
%%% Feature List %%%
%%%%%%%%%%%%%%%%%%%%

\section{Feature List}
\label{sec:rng-features}

The random number generator in \chipname{} generates true random numbers, which means random numbers generated from a physical process, rather than by means of an algorithm. No number generated within the specified range is more or less likely to appear than any other number.

%%%%%%%%%%%%%%%%%%%%%%%%%%%%%%
%%% Functional Description %%%
%%%%%%%%%%%%%%%%%%%%%%%%%%%%%%

\section{Functional Description}
\label{sec:rng-func-descr}

Every 32-bit value that the system reads from the \hyperref[regdesc:LPSYSREGRNGDATAREG]{LPSYSREG\_RNG\_DATA\_REG} register of the random number generator is a true random number. These true random numbers are generated based on the system’s \textbf{thermal noise}, the \textbf{asynchronous clock mismatch} and the \textbf{BUF\_CHAIN}.

\begin{itemize}
    \item \textbf{Thermal noise} comes from the SAR ADC. Whenever the SAR ADC is enabled, bit streams will be generated and fed into the random number generator through an XOR logic gate as random seeds.
    \item RC\_FAST\_CLK is an \textbf{asynchronous clock} source and it increases the RNG entropy by introducing circuit metastability. Please refer to the \hyperref[others:rng-note-clock-seed]{Note} about the RC\_FAST\_CLK.
    \item The count value generated by the \textbf{BUF\_CHAIN}, implemented using BUF\_CELL, is used as a random seed for the random number generator.
\end{itemize}

\begin{figure}[ht]
    \centering
    \includegraphics[width=1.0\textwidth]{\modulefiles/figures/34-RNG-noise-source-latest}
    \caption{Noise Source}
    \label{fig:rng-noise-source}
\end{figure}

When there is noise coming from the SAR ADC, the random number generator is fed with a 1-bit entropy in one clock cycle of RC\_FAST\_CLK, which is generated from an internal RC oscillator (see Chapter \ref{mod:resclk} \textit{\nameref{mod:resclk}} for details). Thus, it is advisable to read the \hyperref[regdesc:LPSYSREGRNGDATAREG]{LPSYSREG\_RNG\_DATA\_REG} register at a maximum rate of 1 MHz to obtain the maximum entropy.


\section{Programming Procedure}

When using the random number generator, make sure at least either the SAR ADC or RC\_FAST\_CLK is enabled. Otherwise, pseudo-random numbers will be returned.
\begin{itemize}
    \item For SAR ADC, please refer to Chapter \ref{mod:adc} \textit{\nameref{mod:adc}}.
    \item RC\_FAST\_CLK is enabled by setting the \hyperref[fielddesc:PMU_HPSLEEPXPDFOSCCLK]{PMU\_HP\_SLEEP\_XPD\_FOSC\_CLK} bit in the \hyperref[regdesc:PMUHPSLEEPLPCKPOWERREG]{PMU\_HP\_SLEEP\_LP\_CK\_POWER\_REG} register. Please refer to the \hyperref[others:rng-note-clock-seed]{Note} about the RC\_FAST\_CLK.
    \item RC\_FAST\_CLK is gated by setting the \hyperref[fielddesc:LPPERICKENRNG]{LPPERI\_CK\_EN\_RNG} bit in the \hyperref[regdesc:LPPERICLKENREG]{LPPERI\_CLK\_EN\_REG} register.
    \item The BUF\_CHAIN can be enabled by setting the \hyperref[regdesc:RNGCFGREG]{RNG\_CFG\_REG} register’s \hyperref[fielddesc:RNGSAMPLEENABLE]{RNG\_SAMPLE\_ENABLE} bit.
\end{itemize}

\begin{tiplisting}\label{others:rng-note-clock-seed}
     Enabling RC\_FAST\_CLK increases the RNG entropy. However, to ensure maximum entropy, it's recommended to always enable an ADC source as well.
\end{tiplisting}

When using the random number generator, read the \hyperref[regdesc:LPSYSREGRNGDATAREG]{LPSYSREG\_RNG\_DATA\_REG} register multiple times until sufficient random numbers have been generated. Ensure the rate at which the register is read does not exceed the frequencies described in section \ref{sec:rng-func-descr} above.


%%%%%%%%%%%%%%%%%%%%%%%%
%%% Register Summary %%%
%%%%%%%%%%%%%%%%%%%%%%%%

\newpage
\hypertarget{rng-reg-summ}{}
\section{Register Summary}
\label{sec:rng-reg-summ}




The abbreviations given in Column \textbf{Access} are explained in Section \hyperref[glossary-access-types]{\textit{Access Types for Registers}}.

\begin{table}[H]
    \centering
    \setlength{\extrarowheight}{0.3em}
        \begin{tabular}{ | l | l | l | l | }
    \hline
    \rowcolor{lightgray}
    \textbf{Name} & \textbf{Description} & \textbf{Address} & \textbf{Access} \\ \hline
    \hyperref[regdesc:LPSYSREGRNGDATAREG]{LPSYSREG\_RNG\_DATA\_REG} & Random source data & 0x5011\_01A4 & RO \\ \hline
    \hyperref[regdesc:LPPERICLKENREG]{LPPERI\_CLK\_EN\_REG} & Integrated clock gating enable signal of RC\_FAST\_CLK & 0x5012\_0000 & R/W \\ \hline
    \hyperref[regdesc:RNGCFGREG]{RNG\_CFG\_REG} & Configuration register of RNG & 0x5012\_6000 & R/W \\ \hline
    \end{tabular}
\end{table}



%%%%%%%%%%%%%%%%%
%%% Registers %%%
%%%%%%%%%%%%%%%%%

\newpage
\section{Registers}


For how to program reserved fields, please refer to Section \hyperref[programming-reserved-field]{Programming Reserved Register Field}.

\begin{register}{H}{LPSYSREG\_RNG\_DATA\_REG}{0x5011\_01A4}\label{regdesc:LPSYSREGRNGDATAREG}
\regfield{LPSYSREG\_RNG\_DATA}{32}{0}{{0x00000000}}%
\reglabel{Reset}\regnewline%
\begin{regdesc}\begin{reglist}
\label{fielddesc:LPSYSREGRNGDATA}\item [LPSYSREG\_RNG\_DATA] Random number source. (RO)
\end{reglist}\end{regdesc}
\end{register}


\begin{register}{H}{RNG\_CFG\_REG}{0x5012\_6000}\label{regdesc:RNGCFGREG}
\regfield{reserved}{22}{10}{{0x000}}%
\regfield{RNG\_TIMER\_EN}{1}{9}{{0}}%
\regfield{RNG\_TIMER\_PSCALE}{8}{1}{{0x00}}%
\regfield{RNG\_SAMPLE\_ENABLE}{1}{0}{{0}}%
\reglabel{Reset}\regnewline%
\begin{regdesc}\begin{reglist}

\label{fielddesc:RNGSAMPLEENABLE}\item [RNG\_SAMPLE\_ENABLE]
Configures whether to enable the RNG asynchronous clock noise source.\\
1: Enable\\
0: Disable \\(R/W)

\label{fielddesc:RNGTIMERPSCALE}\item [RNG\_TIMER\_PSCALE] Represents the prescale factor of the RNG internal timer.(R/W)

\label{fielddesc:RNGTIMEREN}\item [RNG\_RTC\_TIMER\_EN] Configures whether to enable the RNG internal timer. \\
1: Enable and use the counter value of this timer when generating random numbers. \\
0: Disable the timer. \\
(R/W)

\end{reglist}\end{regdesc}
\end{register}
\end{document}

